% !TeX root = ../../Book3.tex
% 纲要标题：### 9.2 高度定理
% 纲要位置：工作记录/计划纲要.md，第 2085 行。
% 纲要细目（写作范围）：
% * Krull 主理想定理
% * 广义高度定理
% * 生成元数与余维约束

\section{高度定理}
\label{b3:sec:0902}

一个方程不应一次切掉任意多层素理想链，这个直觉在 Noether 环中由主理想定理保证。证明需要小心区分极小素理想与任意包含该方程的素理想。本节先用准素性和 Nakayama 引理建立单个方程的结论，再归纳处理多个方程；不使用尚未讲授的 Hilbert–Samuel 理论或 Artin–Rees 引理。

\subsection{Krull 主理想定理}
\label{b3:subsec:0902:01}

\begin{lemma}\label{b3:lem:principal-primary-local-domain}
设 \((R,\mathfrak m)\) 为 Noether 局部整环，且
\(\sqrt{aR}=\mathfrak m\)。则 \(\dim R\le1\)。
\end{lemma}
\begin{proof}
若 \(a=0\)，整环性给出 \(\mathfrak m=0\)，故 \(R\) 是域。以下设 \(a\ne0\)。
反设存在非零素理想 \(0\ne\mathfrak p\subsetneq\mathfrak m\)。因
\(\sqrt{aR}=\mathfrak m\)，有 \(a\notin\mathfrak p\)。

环 \(R/aR\) 是零维 Noether 环，所以由\cref{b3:thm:noether-dimension-zero}为 Artin 环。考虑符号幂
\(Q_n=\mathfrak p^nR_{\mathfrak p}\cap R\)。在 \(R_{\mathfrak p}\) 中，\(\mathfrak p^nR_{\mathfrak p}\) 的根为唯一极大理想，故由\cref{b3:prop:maximal-radical-primary}准素；再由\cref{b3:prop:primary-localization}，收缩后的每个 \(Q_n\) 都是 \(\mathfrak p\) 准素理想。它们在 \(R/aR\) 中的像组成降链，因此存在 \(n\ge1\) 使
\[
Q_n+aR=Q_{n+1}+aR.
\]
对 \(x\in Q_n\)，可写 \(x=y+ar\)，其中 \(y\in Q_{n+1}\)。于是 \(ar\in Q_n\)；
因 \(a\notin\sqrt{Q_n}=\mathfrak p\)，准素性迫使 \(r\in Q_n\)。所以
\[
Q_n=Q_{n+1}+aQ_n.
\]
有限生成模 \(Q_n/Q_{n+1}\) 因而等于自身的 \(a\) 倍。由 \(a\in\mathfrak m\) 及
\cref{b3:thm:nakayama}，它为零，故 \(Q_n=Q_{n+1}\)。

在 \(\mathfrak p\) 处局部化得到
\(\mathfrak p^nR_{\mathfrak p}=\mathfrak p^{n+1}R_{\mathfrak p}\)。
再对 Noether 局部环 \(R_{\mathfrak p}\) 上的有限模
\(\mathfrak p^nR_{\mathfrak p}\) 使用 Nakayama 引理，得到它为零。但
\(R_{\mathfrak p}\) 是整环，且非零的 \(\mathfrak p\) 在其中仍非零，其正整数次幂不可能为零，矛盾。于是除零理想外，没有严格小于 \(\mathfrak m\) 的素理想。
\end{proof}

\begin{theorem}[Krull 主理想定理]\label{b3:thm:principal-ideal-theorem}
设 \(R\) 为交换 Noether 环，\(a\in R\)。每个在 \((a)\) 上极小的素理想 \(\mathfrak p\) 都满足
\[
\operatorname{ht}\mathfrak p\le1.
\]
\end{theorem}
\begin{proof}
在 \(\mathfrak p\) 处局部化后，只存在一个包含 \(a\) 的素理想，即唯一极大理想；
因此可设 \(R\) 为局部环且 \(\sqrt{aR}=\mathfrak m\)。若存在长度至少二的素链，可截取并在末端补上 \(\mathfrak m\)，得到
\(\mathfrak q_0\subsetneq\mathfrak q_1\subsetneq\mathfrak m\)。
在局部整环 \(R/\mathfrak q_0\) 中，\(a\) 的像的根仍为极大理想的像，上一引理却排除了非零的中间素理想
\(\mathfrak q_1/\mathfrak q_0\)，矛盾。因此局部维数至多为一，即所求高度界。
\end{proof}

这称为\term[Krull-zhulixiang-dingli]{Krull 主理想定理}[Krull's principal ideal theorem]。
在整环中，若 \(a\) 非零且非单位，其上极小素理想非零，故高度恰好为一。
“极小”不可删除：例如 \((x)\subseteq(x,y)\) 于 \(k[x,y]\)，后者包含主理想
\((x)\)，却有链 \(0\subsetneq(x)\subsetneq(x,y)\)，高度至少为二。

\subsection{广义高度定理}
\label{b3:subsec:0902:02}

\begin{theorem}[广义高度定理]\label{b3:thm:height-theorem}
若 \(R\) 为交换 Noether 环，\(I=(a_1,\ldots,a_r)\subsetneq R\)，其中 \(r\ge0\)，则每个在 \(I\) 上极小的素理想 \(\mathfrak p\) 都满足
\[
\operatorname{ht}\mathfrak p\le r.
\]
\end{theorem}
\begin{proof}
对 \(r\) 归纳。\(r=0\) 时，\(\mathfrak p\) 是极小素理想，高度为零；\(r=1\) 已由主理想定理证明。一般情形在 \(\mathfrak p\) 处局部化，可设
\((R,\mathfrak m)\) 局部且 \(\sqrt I=\mathfrak m\)。

取任意链
\[
\mathfrak p_0\subsetneq\cdots\subsetneq\mathfrak p_d=\mathfrak m.
\]
若 \(d=0\) 无须证明。因 \(I\not\subseteq\mathfrak p_{d-1}\)，重排生成元后可设
\(a_1\notin\mathfrak p_{d-1}\)。令 \(\mathfrak q_d=\mathfrak m\)，从
\(i=d-1\) 向下逐次在 \(\mathfrak q_{i+1}\) 内选一个在
\(\mathfrak p_i+(a_1)\) 上极小的素理想 \(\mathfrak q_i\)。这样的选择可以先在
\(\mathfrak q_{i+1}\) 处局部化、取一个极小素理想，再收缩；任何更小的候选也在
\(\mathfrak q_{i+1}\) 内，故它在原环中仍然极小。

对 \(i\le d-2\)，若 \(\mathfrak q_i=\mathfrak q_{i+1}\)，则
\[
\mathfrak p_i\subsetneq\mathfrak p_{i+1}\subsetneq\mathfrak q_i.
\]
第二个包含严格，因为 \(a_1\in\mathfrak q_i\) 而
\(a_1\notin\mathfrak p_{i+1}\)。但在 \(R/\mathfrak p_i\) 中，
\(\mathfrak q_i/\mathfrak p_i\) 在一个主理想上极小，主理想定理禁止它具有这样的长度二链。故
\(\mathfrak q_0\subsetneq\cdots\subsetneq\mathfrak q_{d-1}\) 严格。

这是一条包含 \(a_1\) 的长度 \(d-1\) 的链。在局部商环 \(R/(a_1)\) 中，
其极大理想是剩余 \(r-1\) 个生成元的根，由归纳假设高度至多 \(r-1\)。
所以 \(d-1\le r-1\)，即 \(d\le r\)。原链任意，结论成立。
\end{proof}

这一定理也称\term[Krull-gaodu-dingli]{Krull 高度定理}[Krull's height theorem]。
证明没有将
\(\operatorname{ht}\mathfrak p\) 擅自拆成两个高度之和；一般环上的素链未必具有允许这种拆分的长度性质。Matsumura 的《Commutative Ring Theory》给出相同高度结论\cite[Theorem 13.5, pp. 100--101]{Matsumura1989CommutativeRingTheory}，但其证明次序先用 Samuel 函数；本卷采用上述独立论证，使后面的增长理论能够调用高度定理。

\subsection{生成元数与余维约束}
\label{b3:subsec:0902:03}

若 \((R,\mathfrak m)\) 为 Noether 局部环，\(\mathfrak m\) 有有限个生成元，故
\cref{b3:thm:height-theorem}说明 \(\dim R<\infty\)。更一般地，若
\(\sqrt{(a_1,\ldots,a_r)}=\mathfrak m\)，就有 \(\dim R\le r\)。
这限制的是使零点集中到闭点所需的方程数；不表示任一主理想的任一零点都只有余维一。

\begin{lemma}\label{b3:lem:dimension-cutting}
设 \((R,\mathfrak m)\ne0\) 为交换 Noether 局部环，\(d=\dim R\)。存在
\(x_1,\ldots,x_d\in\mathfrak m\)，使
\[
\sqrt{(x_1,\ldots,x_d)}=\mathfrak m.
\]
\(d=0\) 时取空组。
\end{lemma}
\begin{proof}
对有限整数 \(d\) 归纳。\(d=0\) 时唯一素理想为 \(\mathfrak m\)，所以幂零根为
\(\mathfrak m\)，空组满足要求。若 \(d>0\)，环的极小素理想只有有限多个，且都严格小于
\(\mathfrak m\)；否则 \(\mathfrak m\) 极小便排除了其他素理想，维数将为零。
由\cref{b3:lem:finite-prime-avoidance}选
\(x_1\in\mathfrak m\) 避开全部极小素理想。

任一包含 \(x_1\) 的素链，其首项不是极小素理想，可以在下面补上一项。
因此 \(\dim R/(x_1)\le d-1\)。对这个局部商环使用归纳假设，提升其至多
\(d-1\) 个元素，得到一个根为 \(\mathfrak m\)、至多由 \(d\) 个元素生成的理想。
若少于 \(d\) 个，则广义高度定理会给出 \(\dim R<d\)，矛盾。所以总数恰为 \(d\)，完成归纳。
\end{proof}

这个引理与高度定理合起来，将局部维数解释为生成一个 \(\mathfrak m\) 准素理想所需的最少元素个数。\secref{b3:sec:0904}将把这样的元素组命名为参数系。它们生成的理想通常小于 \(\mathfrak m\)，因此参数个数和极大理想的最少生成元数仍需区分。
